\chapter{Kesimpulan dan Saran}\label{ch:summary}
% =============================================================================
% KERANGKA BAB 5 — pola TA rujukan: kesimpulan bernomor menjawab rumusan masalah
% satu-satu dengan angka kunci; saran bernomor untuk penelitian lanjutan.
% Setiap angka di kesimpulan WAJIB sama dengan angka di Bab 4 (rujuk label).
% =============================================================================

    \section{Kesimpulan}\label{sec:conclusion}

        Berdasarkan hasil analisis yang telah dilakukan, kesimpulan yang menjawab rumusan masalah penelitian ini adalah sebagai berikut.
        \begin{enumerate}
            \item % Menjawab rumusan masalah 1 (modifikasi model MCCaR sebagai optimasi berkendala).
                  % Bukti: Pers. 2.21–2.24; verifikasi kasus sintetis (tests/test_risk.py,
                  % tests/test_constraints.py); 360/360 solusi memenuhi kendala bobot
                  % (Tabel \ref{tab:validation feasibility}); kelayakan terhadap batas
                  % bank latih vs batas bank validasi (Tabel \ref{tab:validation true limits}).
                  % SUMBER: results/tables/validation_summary.md
                  % SUMBER: results/tables/validation_true_limits_summary.csv
                  \textit{[diisi]}
            \item % Menjawab rumusan masalah 2 (KMA vs GA vs PSO) pada empat kriteria proposal:
                  % (a) nilai fungsi objektif: Tabel \ref{tab:descriptive stats}, \ref{tab:practical significance};
                  % (b) kelayakan: Tabel \ref{tab:validation feasibility}, \ref{tab:fisher};
                  % (c) konsistensi antarreplikasi: simpangan baku/IQR, Tabel \ref{tab:descriptive stats};
                  % (d) waktu komputasi: Tabel \ref{tab:computation time}.
                  % Sertakan diagnosis penyebab (Subbab \ref{subsec:kma diagnosis});
                  % hasil negatif dinyatakan apa adanya.
                  % SUMBER: results/tables/stats_praktis.csv
                  % SUMBER: results/tables/timing_single_worker_360.csv
                  \textit{[diisi]}
        \end{enumerate}

    \section{Saran}\label{sec:suggestion}

        Berdasarkan hasil yang telah diperoleh, saran untuk penelitian selanjutnya adalah sebagai berikut.
        \begin{enumerate}
            \item % Model skenario dengan mean reversion / perubahan rezim (Subbab \ref{subsec:no mean reversion}).
                  \textit{[diisi]}
            \item % Horizon multiperiode dengan jadwal jatuh tempo dan rollover (Subbab \ref{subsec:one year horizon}).
                  \textit{[diisi]}
            \item % Batas risiko dari beberapa bank atau bank berukuran besar (Subbab \ref{subsec:single bank limits}).
                  \textit{[diisi]}
            \item % Penyetelan parameter algoritma pada anggaran setara (D11 Opsi B sebagai penelitian lanjut;
                  % arah dari ablasi Tahap 18, Subbab \ref{subsec:kma diagnosis}).
                  \textit{[diisi]}
            \item % Studi drift dengan common random numbers dan beberapa bank per skenario (audit T4).
                  \textit{[diisi]}
        \end{enumerate}
